% !TeX program = xelatex
% =============================================================
% 模板 05：IEEE Conference 会议论文（IEEEtran 官方文档类）
% 编译： scripts\build.ps1 en-ieee
% 版式： 双栏，10 pt，官方 IEEEtran.cls（TeX Live 自带，非自制）
% 说明： IEEEtran 默认面向 pdfLaTeX；本工作台统一用 XeLaTeX 编译，
%        若需与官方 Word/PDF 模板逐像素一致，可改用 pdflatex 编译。
% =============================================================
\documentclass[conference]{IEEEtran}

\usepackage{cite}
\usepackage{amsmath}
\usepackage{fontspec}
% IEEEtran 在 XeLaTeX 下不会自动加载 Times，这里显式换成 Times 的度量兼容克隆
\setmainfont{TeX Gyre Termes}
\usepackage{unicode-math}
\setmathfont{TeX Gyre Termes Math}
\usepackage{amsfonts}
\usepackage{graphicx}
\usepackage{booktabs}
\usepackage{textcomp}
\usepackage{xcolor}
\usepackage[hidelinks]{hyperref}

\begin{document}

% ------------------------- 标题区（改这里） -------------------------
\title{Placeholder Title: A Sample IEEE Conference Paper for Template Testing}

\author{
  \IEEEauthorblockN{San Zhang\IEEEauthorrefmark{1} and Si Li\IEEEauthorrefmark{2}}
  \IEEEauthorblockA{\IEEEauthorrefmark{1}School of Electrical Engineering, XX University\\
  Beijing, China \\
  Email: zhangsan@example.edu.cn}
  \IEEEauthorblockA{\IEEEauthorrefmark{2}School of Automation, XX University\\
  Beijing, China \\
  Email: lisi@example.edu.cn}
}

\maketitle

% ------------------------- 摘要 -------------------------
\begin{abstract}
This paper is a placeholder for an IEEE conference submission. It demonstrates the
title block, abstract, a numbered equation, a table, and a bibliography entry so that
the layout can be checked before real content is written. Replace every sentence here
with your own 150--250 word abstract.
\end{abstract}

\begin{IEEEkeywords}
placeholder, template, IEEEtran, conference paper
\end{IEEEkeywords}

\section{Introduction}
\IEEEPARstart{T}{his} is the introduction placeholder. A typical IEEE conference paper
states the problem, reviews related work~\cite{lecun2015deep}, and lists contributions
in a short bullet list at the end of this section.

\section{Method}
Let $x_i$ denote the $i$-th input feature and $w_i$ its weight. The weighted sum is
\begin{equation}
  y=\sum_{i=1}^{n} w_i x_i + b,
  \label{eq:linear}
\end{equation}
where $b$ is the bias term. Equation~\eqref{eq:linear} is the equation placeholder.

\section{Experimental Setup}
The parameter settings used in the placeholder experiment are listed in
Table~\ref{tab:params}.

\begin{table}[h]
  \renewcommand{\arraystretch}{1.2}
  \caption{Placeholder Experimental Parameters}
  \label{tab:params}
  \centering
  \begin{tabular}{lcc}
    \toprule
    Parameter & Symbol & Value \\
    \midrule
    Learning rate & $\eta$ & 0.01 \\
    Batch size    & $B$    & 32   \\
    Epochs        & $E$    & 100  \\
    \bottomrule
  \end{tabular}
\end{table}

\section{Results}
With the settings in Table~\ref{tab:params}, the placeholder model reaches 92.4\%
accuracy on the validation split, which is consistent with the trend reported
in~\cite{he2016resnet}.

\section{Conclusion}
This template compiles end-to-end with XeLaTeX and BibTeX. Replace the placeholder
text, figures, and references with your own content before submission.

\bibliographystyle{IEEEtran}
\bibliography{refs}

\end{document}
